\documentclass[11pt]{amsart}
\usepackage[margin=1.1in]{geometry}
\usepackage{amsmath,amssymb,amsthm}
\usepackage{booktabs}
\usepackage[colorlinks=true,linkcolor=blue,citecolor=blue,urlcolor=blue]{hyperref}

\newtheorem{theorem}{Theorem}
\newtheorem{proposition}[theorem]{Proposition}
\newtheorem{lemma}[theorem]{Lemma}
\newtheorem{corollary}[theorem]{Corollary}
\theoremstyle{definition}
\newtheorem{definition}[theorem]{Definition}
\theoremstyle{remark}
\newtheorem{remark}[theorem]{Remark}

\newcommand{\lbl}[1]{\textup{\textsc{#1}}}
\newcommand{\R}{\mathcal{R}}
\newcommand{\DHL}{\mathrm{DHL}}

\title[Verifying the variational step of the 2026 gap bounds]{An independent exact verification of the variational step behind the 2026 bounds $H_1\le240$ and $H_1\le236$, with a remark on the caps of a $212$ draft}
\author{Jerry Kim}
\thanks{Prepared with the assistance of Claude (Anthropic). All programs, certificates and Lean sources are in the repository \url{https://github.com/jerryckim16/twin-primes}; a living companion paper with full details is at \url{https://jerryckim16.github.io/twin-primes/}.}
\date{Version 0.1 (draft), 13 November 2026}

\begin{document}

\begin{abstract}
Several 2026 preprints claim improvements of the bound $H_1\le246$ of
Polymath8b for $H_1=\liminf(p_{n+1}-p_n)$. Each combines new equidistribution
estimates for primes in arithmetic progressions with a \emph{variational step}:
a test function $F$ on a capped region of $\mathbb R^k$ whose sieve ratio
exceeds a threshold. We recompute the variational step of three of these
claims in exact rational arithmetic, with independent programs, and do not
address their analytic input. (i) For the support and parameters of Stadlmann
($k=49$, claimed $H_1\le240$) we certify the capped ratio $4.0120389\ldots>4$
with a test function of our own, and $4.0124198\ldots$ after optimising it on
the capped region; two independent programs agree to the last digit.
(ii) For the support shared by Kintali and by Song and Yue ($k=48$, claimed
$H_1\le236$) we recompute the published witness of Song and Yue: its cap-free
mass and its capped numerator equal the paper's reference integers ($655$ and
$715$ digits) exactly, so their criterion $48J/I_0>2001/2000$ holds, with
$48J/I_0=1.000664921591\ldots$; the authors list the regeneration of the
capped numerator as pending. The same computation reproduces Kintali's
claimed capped and cap-free values to every printed digit, and a test
function of our own gives $4.0060814\ldots>4$ on the same support. The capped
numerator has so far been recomputed by one program; a second is running.
(iii) For the draft of Charton et al.\ ($k=45$, claimed $H_1\le212$) the
criterion holds for the caps of its Table~3 ($4.0027351\ldots>4$, two
programs), whereas for the different caps stated in its Lemma~9.1 no
polynomial in an $846$-dimensional space of the size the authors use reaches
$3.99$ (ball arithmetic). This locates an inconsistency in the draft; it is
not a refutation. None of this verifies the new equidistribution estimates,
on which every one of these bounds rests, and none of it bears on the twin
prime conjecture, which remains open.
\end{abstract}

\maketitle

\section{Introduction}

\subsection*{The reduction}
For an admissible $k$-tuple $\mathcal H=\{h_1<\dots<h_k\}$, let $\DHL[k,2]$
denote the statement that infinitely many translates $n+\mathcal H$ contain at
least two primes. If $\DHL[k,2]$ holds and $\mathcal H$ is an admissible
$k$-tuple of diameter $D$, then $H_1\le D$. Maynard's method \cite{May15}
proves $\DHL[k,2]$ from a level of distribution and a symmetric test function
$F$ whose ratio $kJ(F)/I(F)$ exceeds a threshold; Polymath8b \cite{Pol14b}
enlarged the region on which $F$ may live and obtained $H_1\le246$ with
$k=50$. The 2026 preprints listed below enlarge the region further, to
\emph{capped} supports, and pay for this with new equidistribution estimates
for the moduli that arise.

\subsection*{The claims examined}
\begin{center}
\begin{tabular}{@{}llll@{}}
\toprule
claim & $k$ & source & what we check\\
\midrule
$H_1\le240$ & $49$ & Stadlmann \cite{Sta26} & variational criterion\\
$H_1\le236$ & $48$ & Kintali \cite{Kin26}; Song--Yue \cite{SY26} & variational criterion; published witness\\
$H_1\le212$ & $45$ & Charton et al.\ \cite{Cha26} & variational criterion; caps\\
\bottomrule
\end{tabular}
\end{center}
Each claim has two parts: an \emph{analytic} part (an equidistribution
theorem for the relevant family of moduli) and a \emph{variational} part (a
function $F$ on the capped support with ratio above the threshold). We check
only the variational part, and the combinatorial fact that an admissible
tuple of the stated diameter exists. Throughout we label statements
\lbl{Computed} (exact certificate and a script that reproduces it),
\lbl{Proved here}, \lbl{Lean-checked}, \lbl{Literature} (read in the source),
or \lbl{Claimed} (asserted in a preprint and not verified by us).

\subsection*{Results in brief}
Ratios in the Polymath scaling, where the threshold is $4$:
\begin{center}\small
\begin{tabular}{@{}lll@{}}
\toprule
claim & variational ratio & status here\\
\midrule
$240$ & $4.0120389$; optimised on $T$: $4.0124198$ & \lbl{Computed}, two programs\\
$236$, witness of \cite{SY26} & $48J/I_0=1.00066492>2001/2000$ & \lbl{Computed}, Thm.~\ref{thm:sy}\\
$236$, our own $F$ & $4.0060814$ & \lbl{Computed}, one program\\
$212$, Table~3 caps & $4.0027351$ & \lbl{Computed}, two programs\\
$212$, Lemma~9.1 caps & $<3.99$ for all $F$ in an $846$-dim.\ space & \lbl{Computed}, ball arithm.\\
\bottomrule
\end{tabular}
\end{center}

\subsection*{What this note does not do}
It does not verify the equidistribution estimates of \cite{Sta26},
\cite{Kin26}, \cite{SY26} or \cite{Cha26}; all four bounds remain conditional
on them, and the established bound is still $246$. Song and Yue
\cite[Remark~4.23]{SY26} point to a missing hypothesis in one step of the
proof of \cite[Lemma~5]{Sta26}; we read that step and comment on it in
Remark~\ref{rem:lemma5}, without resolving it. Nothing here concerns the
parity barrier or the twin prime conjecture.

\section{The variational problem on a capped support}

\subsection*{Supports}
We use the Polymath scaling $u=4t$, in which the criterion reads $kJ/I>4$. Let
$\R_k=\{u\in[0,1]^k:\ u_1+\dots+u_k\le1\}$ and write $L\R_k$ for its dilate.
The one-band supports of \cite{Sta26} (\lbl{Literature}: Definitions~1
and~5 and Proposition~1 of \cite{Sta26}) are described by a length $L$, a marginal radius
$\tau$, a threshold $g$ and caps $C_1\le C_2\le\dots$; with
$I(u)=\{i:u_i>g\}$, $m(u)=|I(u)|$ and $s(u)=\sum_{i\in I(u)}u_i$,
\[
T=\{u\in L\R_k:\ s(u)\le C_{m(u)}\}.
\]
For a symmetric $F$ supported on $T$,
\[
I_T(F)=\int_TF^2,\qquad
J_T(F)=\int_{\tau\R_{k-1}}\left(\int F(u',u_k)\,du_k\right)^2du',
\]
and the variational criterion is $kJ_T(F)/I_T(F)>4$. In this scaling the
parameters are as follows (\lbl{Literature}).
\begin{center}\small
\begin{tabular}{@{}lllll@{}}
\toprule
 & $k$ & $(L,\tau)$ & $g$ & caps $C_1,C_2,\dots$\\
\midrule
\cite{Sta26} & $49$ & $(1.042,\,0.982)$ & $0.112$ & $0.6,\,0.6,\,0.68,\,0.68,\dots$\\
\cite{Kin26}, \cite{SY26} & $48$ & $(1.042,\,0.982)$ & $91/1250$ & $1899/2500$ for all $m$\\
\cite{Cha26}, Tab.~3 & $45$ & $(53/50,\,249/250)$ & $41/625$ & $4B_{1,m}$ (see below)\\
\cite{Cha26}, L.~9.1 & $45$ & $(53/50,\,249/250)$ & $41/625$ & $0.62,\,0.62,\,0.68,\,0.68,\dots$\\
\bottomrule
\end{tabular}
\end{center}
For \cite{Cha26}, Table~3 lists
$B_{1,m}=\frac1{5000}(777,794,875,917,953,983,1016,1042,1063,1081)$ for
$m\le10$ and $B_{1,m}=1081/5000$ for $m\ge11$, and $C_m=4B_{1,m}$; Lemma~9.1
states $C_1=C_2=31/50$ and $C_m=17/25$ for $m\ge3$. All four cap sets satisfy
$C_{m+1}\le C_m+g$ (\lbl{Lean-checked} for the caps of \cite{Cha26} and the
counting lemmas below).

\subsection*{Exact moments}
All integrals reduce to moments over simplices.

\begin{lemma}[\lbl{Proved here}; elementary]\label{lem:egf}
Let $P_j=\sum_iu_i^j$, $X_1,\dots,X_k$ independent exponential variables of
mean $1$, and $E_k(b)=\mathbb E\prod_jP_j(X)^{b_j}$. Then, with
$w(b)=\sum_jjb_j$,
\[
\int_{\R_k}(1-P_1)^A\prod_jP_j^{b_j}\,du=\frac{A!}{(k+A+w(b))!}\,E_k(b),
\qquad
\sum_bE_k(b)\frac{s^b}{b!}=\left(\sum_c w(c)!\,\frac{s^c}{c!}\right)^{k}.
\]
\end{lemma}

\begin{proof}
The Dirichlet integral gives $\int_{\R_k}(1-P_1)^Au^\alpha\,du=A!\,\alpha!/(k+A+|\alpha|)!$
and $\alpha!=\mathbb EX^\alpha$; sum over the monomials of $\prod_jP_j^{b_j}$,
all of degree $w(b)$. The generating function follows from independence.
\end{proof}

\subsection*{The fibres of a capped support}
Let $R=L-|u'|$ on the fibre over $u'\in\mathbb R^{k-1}_{\ge0}$, and call $u'$
\emph{admissible} if $s(u')\le C_{m(u')}$.

\begin{proposition}[\lbl{Proved here}; elementary]\label{prop:fibre}
If $C_{m+1}\le C_m+g$ for all $m$, then the fibre $\{t\ge0:(u',t)\in T\}$ is
empty when $u'$ is not admissible and equals $[0,U(u')]$ otherwise, where
\[
U=\min\bigl(R,\ \max(g,\ C_{m(u')+1}-s(u'))\bigr).
\]
Consequently, with $Q(u',x)=\int_0^xP(u',t)\,dt$ for a polynomial $P$,
\[
J_T(P\,1_T)=J(P)-\int_{\substack{|u'|\le\tau\\ u'\text{ not adm.}}}Q(u',R)^2
-\int_{\substack{|u'|\le\tau\\ u'\text{ adm.},\ U<R}}\bigl(Q(u',R)^2-Q(u',U)^2\bigr),
\]
and $I_T(P\,1_T)=I(P)-\int_{L\R_k\setminus T}P^2$, where $I,J$ are the
cap-free functionals.
\end{proposition}

\begin{proof}
If $t\le g$ the large set of $(u',t)$ is $I(u')$, so $(u',t)\in T$ if and only
if $u'$ is admissible and $t\le R$. If $t>g$ the large set gains the last
coordinate and the condition is $s(u')+t\le C_{m(u')+1}$; with
$C_{m+1}\le C_m+g$ this forces $s(u')<C_{m(u')}$, so a non-admissible $u'$ has
an empty fibre, and for admissible $u'$ the fibre is
$[0,\min(R,g)]\cup(g,\min(R,C_{m+1}-s)]=[0,U]$. The marginal of $P\,1_T$ is
$Q(u',U)$ on admissible $u'$ and $0$ otherwise; subtract from the cap-free
marginal $Q(u',R)$.
\end{proof}

By symmetry each correction is a finite signed sum, over the size $m$ of the
large set and over $j\ge0$ further coordinates forced above $g$ (inclusion and
exclusion for ``all other coordinates are at most $g$''), of integrals of
polynomials over polygons in two radial variables; so every term is an exact
rational number. The ranges of $m$ and $j$ that occur are short and were
checked in Lean (e.g.\ $m\le9$, $j\le3$ for \cite{Sta26}; $m\le10$, $j\le3$
for \cite{Kin26}; $m\le16$, $j\le6$ for \cite{Cha26}).

\section{Method}

\subsection*{Two independent exact evaluators}
Given a polynomial $P$ with rational coefficients (we use products of
$(L-P_1)^a$ and even power sums), both programs return $I_T(P\,1_T)$ and
$J_T(P\,1_T)$ as exact rationals.
\begin{enumerate}
\item \texttt{capped\_exact.py} computes the cap-free values with
Lemma~\ref{lem:egf} and subtracts the corrections of
Proposition~\ref{prop:fibre}, piece by piece.
\item \texttt{capped\_check.py} was written separately from the definitions
only, without access to the first program. It sums directly over the exact
large set, with exponential-moment generating functions, and has no cap-free
term.
\end{enumerate}
Both were tested independently: a Monte Carlo comparison at $k=12$ with caps
that remove almost half of $I$ and $J$; exact agreement of
\texttt{capped\_check.py} with brute-force polytope triangulation and symbolic
integration for $k=2,\dots,5$; and exact agreement of the two programs on every
production input reported as ``two programs'' below. A value counts as
certified when both programs return the same two rationals.

\subsection*{Optimising on the capped support, and ceilings}
Both functionals are quadratic forms, $I_T=c^\top G_Ic$ and $J_T=c^\top G_Jc$,
in the coefficients $c$ of $P$. Running the chain of exact linear maps of
\texttt{capped\_check.py} backwards gives the Gram matrices at about the cost
of one evaluation (\texttt{capped\_gram.py}; \lbl{Proved here}, elementary).
We use them in two ways. First, to \emph{find} a good vector (the top
generalised eigenvector), which is then rounded and re-evaluated exactly by
the evaluators; only that exact value is stated as a certificate. Second, to
prove a \emph{ceiling}: if $\sigma G_I-kG_J$ is positive definite, no $P$ in
the space reaches the ratio $\sigma$. Positive definiteness is certified by
recursive Schur complements in ball arithmetic, with the entrywise error of
$G_J$ included in the balls. The weights behind $G_I$ and $G_J$ were checked
against \texttt{capped\_check.py} exactly at a uniformly random integer
vector, which by the Schwartz--Zippel lemma rules out a discrepancy between
the two quadratic forms except with probability below $2^{-31}$ for each.

\section{Results}

\subsection*{The support of Stadlmann ($k=49$, claimed $H_1\le240$)}

\begin{theorem}[\lbl{Computed}; exact, two independent programs]\label{thm:240}
For the parameters of \cite{Sta26} in the table above:
\begin{enumerate}
\item for $P$ the cap-free optimiser of degree $20$ in even power sums
($707$ coefficients),
\[49J_T(P1_T)/I_T(P1_T)=4.0120389035\ldots;\]
\item for $P^\ast$ the top eigenvector of $(49G_J,G_I)$ in the same space,
rounded to $400$ significant bits per coefficient,
\[49J_T(P^\ast1_T)/I_T(P^\ast1_T)=4.0124198475\ldots.\]
\end{enumerate}
\end{theorem}

In both cases the two evaluators agree on $I_T$ and $J_T$ as exact rationals.
Without caps the same space gives $4.016176$, so the caps cost $0.0041$ for the
truncated optimiser and $0.0038$ for the optimised one: optimising on $T$
regains less than a tenth of the loss. Since every point of $T$ has all
coordinates at most $0.6$ and $T\subset L\R_{49}$, a dual-weight bound gives
the ceiling $\sup_FkJ_T/I_T\le4.14$ (\lbl{Computed}, interval arithmetic). So
the margin of the variational step of \cite{Sta26} lies in $(0.0124,\,0.14]$.

\begin{corollary}[conditional; the analytic input is \lbl{Claimed}]
The criterion of \cite[Prop.~1]{Sta26} holds for $k=49$ with Stadlmann's
support, caps and parameters. With the \lbl{Lean-checked} admissible
$49$-tuple of diameter $240$, $H_1\le240$ follows from \cite[Prop.~1]{Sta26}
and his equidistribution theorem, which we have not verified.
\end{corollary}

\subsection*{The support of Kintali and of Song--Yue ($k=48$, claimed $H_1\le236$)}
Both preprints use the same support (the table above, called ``D182'' in
\cite{Kin26}) and, by their descriptions, the same public witness of degree
$25$, originally constructed for $k=50$ and evaluated at $k=48$ without
re-optimisation. Song and Yue print it as $1295$ integer triples
$(a,\alpha,c)$, with
\[
P(z)=\sum c\,\left(\tfrac{26}{25}-s(z)\right)^a m_\alpha(z),\qquad
P_U(t)=P\!\left(\tfrac{26}{25U}t\right),\quad U=\tfrac{521}{2000},
\]
$m_\alpha$ the monomial symmetric polynomials, and they print the reference
integers $I_{\mathrm{ref}}=(98!)^2I_0^\ast$ and
$J_{\mathrm{ref}}=(99!)^2(26!)^2J^\ast$. Here $I_0$ is the \emph{cap-free} mass
on $\{s<U\}$, $J$ the capped numerator, and the stars denote integer
coordinates $x=10^4t$. Their criterion is $48J/I_0>2001/2000$, which implies
the capped criterion because $I_T\le I_0$ (\lbl{Literature},
\cite[\S4.12, App.~A, App.~C]{SY26}). They state that a fresh evaluation
reproduces $I_{\mathrm{ref}}$, while fresh regeneration of $J_{\mathrm{ref}}$
from the table ``remains incomplete''.

We extracted the table from the PDF using font sizes to read the exponents in
the partitions. All $1295$ rows are present, every $\alpha$ is nondecreasing
with even parts, and $\max(a+|\alpha|)=25$, as stated. Since all parts are
even, $P_U$ is exactly a polynomial in $L-P_1$ and the even power sums
$P_2,\dots,P_{24}$ ($1311$ terms); the conversion was checked exactly at random
rational points.

\begin{theorem}[\lbl{Computed}, exact rational arithmetic]\label{thm:sy}
For the witness $P_U$ of \cite{SY26} and the support above:
\begin{enumerate}
\item $(98!)^2\,10^{192}\,5^{150}\,521^{50}\,I_0=I_{\mathrm{ref}}$ exactly
(an integer with $655$ digits);
\item $(99!)^2(26!)^2\,10^{196}\,5^{150}\,521^{50}\,J_T=J_{\mathrm{ref}}$
exactly (an integer with $715$ digits), where $J_T$ is the capped numerator;
\item hence $48J_T/I_0=1.000664921591\ldots>2001/2000$, which is
$4.0026596864\ldots$ in the Polymath scaling; with the capped denominator
$48J_T/I_T=4.0030988047\ldots$, and without caps $48J_0/I_0=4.0036348167\ldots$.
\end{enumerate}
\end{theorem}

The factors $10^{192}5^{150}521^{50}$ and $10^{196}5^{150}521^{50}$ come from the
change to integer coordinates when the polynomial is scaled by $65125^{25}$
(our reading of \cite[Eq.~(71)]{SY26}; $65125=25\cdot2605$), so items (1) and
(2) are equalities of integers, not numerical agreement. Item (1) was computed
by a program written for this purpose, which works in the monomial basis of
the table, integrating each product $m_\alpha m_\beta$ by a sum over partial
matchings of parts, and does not use the power-sum conversion. Item (2) was
computed by \texttt{capped\_check.py} on the converted polynomial ($56$ minutes
on four cores). The capped and cap-free values in (3) agree with those claimed
in \cite{Kin26}, $4.0030988$ and $4.0036348$, to every printed digit.

\emph{Status of item (2).} This item and Proposition~\ref{prop:kin-own}
below are the results of this note that so far rest on one program. The second evaluator, \texttt{capped\_exact.py}, agrees
with it on the cap-free part and is computing the $126$ correction pieces
from a checkpoint (about $5$--$10$ hours on four cores). This draft will be
updated when it finishes. The exact agreement of the first program with the
authors' independently produced $715$-digit integer is itself strong evidence
for both, since a common error would have to produce the same integer.

\begin{proposition}[\lbl{Computed}, exact rational arithmetic, one program]\label{prop:kin-own}
For our own cap-free optimiser of degree $22$ for this $(L,\tau)$ at $k=48$
($1041$ coefficients, even power sums, rounded to $400$ significant bits),
$48J_T(P1_T)/I_T(P1_T)=4.0060814208\ldots>4$.
\end{proposition}

Here the caps cost only $0.00033$ (the cap-free value is $4.0064115$). That our
value exceeds the witness's $4.0031$ is expected, since the witness was
optimised for $k=50$ and transported, while ours is optimised for $k=48$.

\begin{corollary}[conditional; the analytic input is \lbl{Claimed}]
The variational criteria of \cite{Kin26} and \cite{SY26} hold at $k=48$, for
the published witness and for an independent test function. With the
\lbl{Lean-checked} admissible $48$-tuple of diameter $236$, $H_1\le236$
follows from their equidistribution input, which we have not verified.
\end{corollary}

\subsection*{The draft of Charton et al.\ ($k=45$, claimed $H_1\le212$)}

\begin{proposition}[\lbl{Computed}; exact, two independent programs]\label{prop:212}
With the caps of \cite[Table~3]{Cha26} and $P$ our cap-free optimiser of
degree $21$ ($846$ coefficients, the dimension of the space
$\mathcal B_{21}$ of \cite[Thm.~11.1]{Cha26}),
$45J_T(P1_T)/I_T(P1_T)=4.0027351128\ldots>4$.
\end{proposition}

The authors report $4.00438$ with their own optimised function (\lbl{Claimed});
our value is lower, as expected for a function optimised without caps, and
already above $4$. The draft, however, also states caps in its Lemma~9.1,
namely $C_1=C_2=31/50$ and $C_m=17/25$ for $m\ge3$, which are smaller than
those of Table~3 for every $m\ge4$. With these the conclusion changes.

\begin{proposition}[\lbl{Computed}, ball arithmetic]\label{prop:ceiling}
Let $T'$ be the support with the caps of Lemma~9.1 of \cite{Cha26} and $V$ the
$846$-dimensional space above. Then $3.99\,G_I-45\,G_J$ is positive definite on
$V$, so
\[
\frac{45J_{T'}(P1_{T'})}{I_{T'}(P1_{T'})}<3.99\qquad\text{for every nonzero }P\in V.
\]
\end{proposition}

The certificate was obtained twice, at $5000$ and $6500$ bits with different
block splits, and the underlying weights agree exactly with
\texttt{capped\_check.py} at a random integer point. The best $P$ in $V$ gives
about $3.98555$ (floating point, from the Gram matrices). So within a space of
the size the authors use, the caps of Lemma~9.1 cannot give their
$4.00438$, which must belong to the caps of Table~3. A careful reading of the
draft supports this: Table~3 is presented as the parameter set used in both
halves of the proof, and the Lemma~9.1 values appear to be left over from an
earlier choice. We report this as a located inconsistency for the authors to
resolve, not as an error in their result. We have not checked that their space
coincides with ours beyond its dimension, and \texttt{capped\_exact.py} has not
yet re-evaluated the random point used for the weights.

\section{Remarks}

\begin{remark}[on {\cite[Lemma~5]{Sta26}}; \lbl{Literature} reading, not a result]\label{rem:lemma5}
Song and Yue \cite[Remark~4.23]{SY26} say that one step in the proof of
\cite[Lemma~5]{Sta26} lacks a hypothesis. In arXiv v1, Lemma~5 (a Type~I
estimate) takes $f=\alpha\star\beta$ with $\alpha$ an arbitrary
coefficient sequence and $\beta$ smooth. For $\gamma$ in a short range above
$1/2$ the proof applies Lemma~3 (a Type~II estimate) ``with the roles of
$\alpha$ and $\beta$ swapped''; Lemma~3 requires its shorter factor to satisfy
a Siegel--Walfisz condition, and after the swap that factor is $\alpha$, for
which neither Lemma~5 nor the case of \cite[Definition~9]{Sta26} where it is
used assumes this. So the step reads as \cite{SY26} describe it. A natural
repair is to check that the sequences $\alpha$ arising in the decomposition
have the property; we have not done so, and we do not know whether the
conclusion is affected. Theorem~\ref{thm:240} concerns only the variational
step and is unaffected; the \emph{bound} $240$ depends on this input.
\end{remark}

\begin{remark}[what would change the picture]
The variational steps examined here have margins of $0.003$ to $0.012$ above
the threshold. They are therefore robust to rounding but not to a change of
support: Proposition~\ref{prop:ceiling} shows that a modest change in the caps
moves the best attainable ratio below $4$. Any revision of the analytic input
that shrinks the admissible caps should be re-checked against the variational
step, which our programs can do in hours.
\end{remark}

\section{Reproducibility}
All inputs, certificates and logs are in the repository, under
\texttt{code/} and \texttt{code/cap6/}. The evaluators are
\texttt{code/capped\_exact.py} and \texttt{code/capped\_check.py}; the Gram
matrices and ceilings come from \texttt{code/capped\_gram.py}; the Song--Yue
witness and reference integers, as extracted, are
\texttt{code/cap6/songyue236\_witness.json} and
\texttt{code/cap6/songyue236\_appC.json}, with the conversion and comparison
scripts \texttt{code/songyue\_*.py}. Running times above are for four cores
and $16$~GB of memory. The admissible tuples of diameters $212$, $236$, $240$
and $246$ are checked by the Lean~4 kernel (core Lean, no Mathlib, no
\texttt{sorry}; axioms \texttt{propext} and \texttt{Quot.sound} only), as are
the counting lemmas used in Proposition~\ref{prop:fibre}. The companion paper
\cite{Kim26} contains the full derivations and the history of each
computation.

\begin{thebibliography}{99}
\bibitem{Cha26} F.~Charton, L.~Hong, K.~Lau, K.~Ono, G.~Remy, H.~C.~Siu,
A.~A.~Swaminathan, J.~Thorner, Y.~Xie, \emph{A new bound for small gaps
between primes}, preliminary draft, 3 September 2026.
\url{https://primegaps.axiommath.ai/bgp212.pdf}
\bibitem{Kim26} J.~Kim, \emph{Toward the twin prime conjecture: certificates,
limits of the sieve, and a working programme}, living paper, 2026.
\url{https://jerryckim16.github.io/twin-primes/}
\bibitem{Kin26} S.~Kintali, \emph{Improved bounded gaps between primes},
preprint, 1 September 2026.
\url{https://shivakintali.github.io/papers/Prime-Gaps.pdf}
\bibitem{May15} J.~Maynard, \emph{Small gaps between primes}, Ann. of Math.
181 (2015), 383--413.
\bibitem{Pol14b} D.~H.~J.~Polymath, \emph{Variants of the Selberg sieve, and
bounded intervals containing many primes}, Res. Math. Sci. 1 (2014).
\bibitem{SY26} Z.~Song, S.~Yue, \emph{Bounded gaps between primes: an upper
bound of 236}, IACR ePrint 2026/1893, version of 2 September 2026.
\bibitem{Sta26} J.~Stadlmann, \emph{Bounded gaps between primes},
arXiv:2608.31126, 31 August 2026.
\end{thebibliography}

\end{document}
